\section{Related Work}\label{sec:related}
The FJ model, integrating internal opinions with social influence, stands as a pivotal framework in social dynamics research~\cite{FrJo90}. Stability conditions for this model were established in~\cite{ravazzi2014ergodic}, while equilibrium expressed opinion was derived in~\cite{BiKlOr15,das2013debiasing}. Interpretations of the FJ model were further explored in~\cite{BiKlOr15} and~\cite{ghaderi2014opinion}. Research has also delved into adversarial interventions~\cite{chen2021adversarial,gaitonde2020adversarial,tu2023adversaries} and the spread of viral content~\cite{tu2022viral}. Furthermore, the FJ model has been expanded in various dimensions, as discussed in~\cite{jia2015opinion,semonsen2018opinion}, including multidimensional approaches~\cite{friedkin2016network,parsegov2016novel}.

Beyond the foundational aspects, interpretations, and extensions of the FJ model, researchers have developed metrics to quantify disagreement~\cite{musco2018minimizing}, polarization~\cite{matakos2017measuring,musco2018minimizing}, conflict~\cite{chen2018quantifying}, and controversy~\cite{chen2018quantifying}. These metrics offer valuable insights into social dynamics. Various optimization problems based on the FJ model have been addressed, including efforts to minimize overall opinion~\cite{sunOpinionOptimizationDirected2023}, disagreement~\cite{musco2018minimizing}, polarization~\cite{musco2018minimizing, matakos2017measuring}, and conflict~\cite{chen2018quantifying,zhu2022nearly} by adjusting node attributes such as internal opinions or susceptibility to persuasion~\cite{abebe2018opinion, gionis2013opinion,chan2019revisiting}. Network modifications, like edge additions or deletions to reduce polarization or foster consensus, have emerged as a significant area of study~\cite{zhu2021minimizing,musco2018minimizing, chen2018quantifying}.

Scalability and efficiency pose significant challenges in large networks, where conventional methods like matrix inversion are impractical. To overcome this, scalable algorithms utilizing sampling and distributed computing have been introduced~\cite{ravazzi2014ergodic, avena2018two}, facilitating the analysis of networks with millions of nodes. For undirected graphs, Laplacian solvers can approximate equilibrium opinions in near-linear time~\cite{xu2021fast}, while directed graphs rely on forest sampling methods~\cite{sunOpinionOptimizationDirected2023}. Recent work~\cite{neumann2024sublinear} propose sublinear time algorithms to estimate single-node opinions without computing the entire equilibrium vector.
